\begin{table}[htbp!]
\centering
\caption{Concordância na anotação às cegas do lote de validação}
\label{tab:lote_concordancia}
\footnotesize
\ajustartabela{%
\begin{tabular}{|l|c|c|}
\hline
\textbf{Medida} & \textbf{Construção $\times$ A} (n=1.950) & \textbf{A $\times$ B} (n=390) \\
\hline
Decisão (buscar / buscar ou esclarecer / não buscar), kappa & 1,000 & 1,000 \\
Leituras compatíveis nos dois sentidos & 99,8\% & 100,0\% \\
Leitura preferencial idêntica & 99,8\% & 100,0\% \\
Presença dos campos, kappa & 0,999 & 1,000 \\
Valor igual quando ambos preenchem & 100,0\% & 100,0\% \\
Detecção de leituras múltiplas, kappa & 0,967 & 1,000 \\
\hline
\end{tabular}}
\fonte{Elaborada pelos autores. Medidas de \texttt{audit.concordancia}, as mesmas da auditoria das 310 consultas, sobre todas as consultas redigidas (antes do filtro). A: anotação de todas as consultas; B: segunda anotação independente de uma amostra aleatória de 20\%.}
\end{table}
